\begin{abstract}
Small-loss sample selection, the core heuristic of learning with noisy labels, stops training on
examples the model currently fits poorly. We show that in the masked-softmax family of AS-Softmax
this heuristic is a margin. A sample's loss and gradient vanish exactly when its label gap
$g = \max_{j\neq y} p_j - p_y$ satisfies $g \le -\delta$, so a \emph{negative} per-sample margin
$\delta$ rejects the sample, and every selection rule is a choice of margin function. The view
turns the rule's suspicion statistic into the object of design and yields testable predictions.
A batch-relative statistic is invariant to the batch's overall fit, so the share it rejects cannot
follow the noise rate. An absolute statistic (``another class beats the label'') rejects a
closed-form band of gaps that does not depend on the noise rate, and leaves an escape region of
confidently contradicted labels that no margin in the family can reject. We test these predictions
in a pre-registered study: BERT-base on 20~Newsgroups, six synthetic noise conditions, five seeds,
model selection on noisy held-out labels. The batch-relative margin ties the best robust loss at
40\% symmetric noise but loses 3.8 points on clean data and 4.6 under pair-flip noise, as its
noise-blindness predicts. The absolute margin memorizes the fewest corrupted labels in four of the
five noisy conditions (8--10\% under symmetric noise, against 13--29\% for small-loss selection
given the true noise rate) and is the best method under 40\% pair-flip noise (+5.4 points over
cross-entropy, +3.3 over oracle small-loss). It still trails GCE by 0.9--1.3 points on clean and symmetric noise, so
it fails our pre-registered bar for a new method. Instrumented re-runs locate the costs. A rejected
clean sample receives no gradient and stays locked out (23\% of a noise-free training set under the
batch-relative margin). AS-Softmax, the base loss, initially concentrates its gradient on
mislabeled samples and memorizes them faster than cross-entropy. We report the result as a unifying
view with measured consequences, not as a state-of-the-art method.
\end{abstract}
